% -*- coding: utf-8 -*-
% !TEX program = xelatex
\documentclass[plain,noanswer,medmath]{../../template/jnuexam}
\usepackage{../../template/kysx}

\begin{document}


\examtitle{2000}{3}


\loadproblems{1}{2000P1}%载入试卷一


\makepart{填空题}{本题共5小题，每小题3分，满分15分}


\begin{problem}
设$z = f\left(xy,\frac{x}{y} \right) + g\left(\frac{x}{y} \right)$，其中$f,g$均可微，则$\frac{\pdz}{\pdx} =$ \fillin{}．
\end{problem}


\begin{problem}
$\int_1^{+\infty} \frac{\dx}{\e^x + \e^{2 - x}} =$ \fillin{}．
\end{problem}


\begin{problem}
若四阶矩阵$A$与$B$相似，矩阵$A$的特征值为$\frac{1}{2},\frac{1}{3},\frac{1}{4},\frac{1}{5}$，
则行列式$\left|B^{-1}-E\right|=$ \fillin{}．
\end{problem}


\begin{problem}
设随机变量$X$的概率密度为$f(x) = \begin{cases}
  1/3, & x \in[0,1]; \\
  2/9, & x \in[3,6]; \\
    0, & \text{其他}.
\end{cases}$ 若$k$使得$P\{X \ge k\} = \frac{2}{3}$，则$k$的取值范围是 \fillin{}．
\end{problem}


\begin{problem}
假设随机变量$X$在区间$[- 1,2]$上服从均匀分布，随机变量$Y = \begin{cases}
   1, & \text{若}\;X > 0; \\
   0, & \text{若}\;X = 0; \\
  -1, & \text{若}\;X < 0.
\end{cases}$
则方差$D(Y) =$ \fillin{}．
\end{problem}


\makepart{选择题}{本题共5小题，每小题3分，共15分}


\begin{problem}
设对任意的$x$，总有$\varphi (x) \le f(x) \le g(x)$，且$\lim_{x \to \infty} \left[g(x) - \varphi(x) \right] = 0$，
则$\lim_{x \to \infty} f(x)$ \pickout{}
\begin{abcd}
\item 存在且一定等于零．
\item 存在但不一定等于零．
\item 一定不存在．
\item 不一定存在．
\end{abcd}
\end{problem}


\begin{problem}
设函数$f(x)$在点$x = a$处可导，则函数$\left| f(x) \right|$在点$x = a$处不可导的充分条件是\pickout{}
\begin{abcd}
\item $f(a) = 0$且$f'(a) = 0$．
\item $f(a) = 0$且$f'(a) \ne 0$．
\item $f(a) > 0$且$f'(a) > 0$．
\item $f(a) < 0$且$f'(a) < 0$．
\end{abcd}
\end{problem}


\begin{problem}
设$\alpha_1,\alpha_2,\alpha_3$是四元非齐次线性方程组$AX = b$的三个解向量，
且秩$r(A) = 3$，$\alpha_1 = \left(1,2,3,4 \right)^T$，$\alpha_2 + \alpha_3 = \left(0,1,2,3 \right)^T$，
$c$表示任意常数，则线性方程组$AX = b$的通解$X =$\pickout{}
\begin{abcd}
\item $\begin{pmatrix}
  1 \\
  2 \\
  3 \\
  4
\end{pmatrix} + c\begin{pmatrix}
  1 \\
  1 \\
  1 \\
  1
\end{pmatrix}$．
\item $\begin{pmatrix}
  1 \\
  2 \\
  3 \\
  4
\end{pmatrix} + c\begin{pmatrix}
  0 \\
  1 \\
  2 \\
  3
\end{pmatrix}$．
\item $\begin{pmatrix}
  1 \\
  2 \\
  3 \\
  4
\end{pmatrix} + c\begin{pmatrix}
  2 \\
  3 \\
  4 \\
  5
\end{pmatrix}$．
\item $\begin{pmatrix}
  1 \\
  2 \\
  3 \\
  4
\end{pmatrix} + c\begin{pmatrix}
  3 \\
  4 \\
  5 \\
  6
\end{pmatrix}$．
\end{abcd}
\end{problem}


\begin{problem}
设$A$为$n$阶实矩阵，$A^T$是$A$的转置矩阵，则对于线性方程组(Ⅰ)：$AX = 0$和(Ⅱ)：$A^T AX = 0$，必有\pickout{}
\begin{abcd}
\item (Ⅱ)的解是(Ⅰ)的解，(Ⅰ)的解也是(Ⅱ)的解．
\item (Ⅱ)的解是(Ⅰ)的解，但(Ⅰ)的解不是(Ⅱ)的解．．
\item (Ⅰ)的解不是(Ⅱ)的解，(Ⅱ)的解也不是(Ⅰ)的解．．
\item (Ⅰ)的解是(Ⅱ)的解，但(Ⅱ)的解不是(Ⅰ)的解．
\end{abcd}
\end{problem}


\begin{problem}
在电炉上安装了$4$个温控器，其显示温度的误差是随机的．
在使用过程中，只要有两个温控器显示的温度不低于临界温度$t_0$，电炉就断电．
以$E$表示事件“电炉断电”，而$T_{(1)} \le T_{(2)} \le T_{(3)} \le T_{(4)}$为$4$个温控器显示的按递增顺序排列的温度值，
则事件$E$等于事件\pickout{}
\begin{abcd}
\item $\left\{T_{(1)} \ge t_0 \right\}$．
\item $\left\{T_{(2)} \ge t_0 \right\}$．
\item $\left\{T_{(3)} \ge t_0 \right\}$．
\item $\left\{T_{(4)} \ge t_0 \right\}$．
\end{abcd}
\end{problem}


\makepart{}{本题满分6分}

\begin{problem}
求微分方程$y'' - 2y' - \e^{2x} = 0$满足条件$y(0) = 0,y'(0) = 1$．
\end{problem}


\makepart{}{本题满分6分}

\begin{problem}
计算二重积分$\iint_D \frac{\sqrt{x^2 + y^2}}{\sqrt{4a^2 - x^2 - y^2}}\d\sigma$，
其中$D$是由曲线$y = - a + \sqrt{a^2 - x^2} (a > 0)$ 和直线$y = - x$围成的区域．
\end{problem}


\makepart{}{本题满分6分}

\begin{problem}
假设某企业在两个相互分割的市场上出售同一种产品，两个市场的需求函数分别是
$$P_1 = 18 - Q_1,P_2 = 12 - Q_2,$$
其中$P_1$和$P_2$分别表示该产品在两个市场的价格(单位：万元/吨)，
$Q_1$和$Q_2$分别表示该产品在两个市场的销售量(即需求量，单位：吨)，
并且该企业生产这种产品的总成本函数是$C = 2Q + 5$，其中$Q$表示该产品在两个市场的销售总量，即$Q = Q_1 + Q_2$．
\begin{enumerate}
\item 如果该企业实行价格差别策略，试确定两个市场上该产品的销售量和价格，使该企业获得最大利润；
\item 如果该企业实行价格无差别策略，试确定两个市场上该产品的销售量及其统一的价格，
      使该企业的总利润最大化；并比较两种价格策略下的总利润大小．
\end{enumerate}

\end{problem}


\makepart{}{本题满分7分}

\begin{problem}
求函数$y = (x - 1)\e^{\frac{\pi}{2} + \arctan x}$的单调区间和极值，并求该函数图形的渐近线．
\end{problem}


\makepart{}{本题满分6分}

\begin{problem}
设$I_n = \int_0^{\frac{\pi}{4}} \sin^n xcosx\dx ,n = 0,1,2, \cdots $，求$\sum_{n = 0}^{\infty}I_n$．
\end{problem}


\makepart{}{本题满分6分}%【同试卷一第九题】

\useproblem{1}{9}{0}


\makepart{}{本题满分8分}

\begin{problem}
设向量组$\alpha_1 = (a,2,10)^T$, $\alpha_2 = (- 2,1,5)^T$, $\alpha_3 = (- 1,1,4)^T,\beta = (1,b,c)^T$．
试问$a,b,c$满足什么条件时，
\begin{enumerate}
\item $\beta$可由$\alpha_1,\alpha_2,\alpha_3$线性表示，且表示唯一？
\item $\beta$不能由$\alpha_1,\alpha_2,\alpha_3$线性表示？
\item $\beta$可由$\alpha_1,\alpha_2,\alpha_3$线性表示，但表示不唯一？并求出一般表达式．
\end{enumerate}
\end{problem}


\makepart{}{本题满分9分}

\begin{problem}
设有$n$元实二次型
$$f(x_1,x_2, \cdots ,x_n) = (x_1 + a_1x_2)^2 + (x_2 + a_2x_3)^2 + \cdots + (x_{n - 1} + a_{n - 1}x_n)^2 + (x_n + a_n x_1)^2,$$
其中$a_i = (i = 1,2, \cdots ,n)$为实数．试问：当$a_1,a_2, \cdots ,a_n$满足何种条件时，
二次型$f(x_1,x_2, \cdots ,x_n)$ 为正定二次型？
\end{problem}


\makepart{}{本题满分8分}

\begin{problem}
假设0.50，1.25，0.80，2.00是来自总体$X$的简单随机样本值．已知$Y = \ln X$服从正态分布$N(\mu ,1)$．
\begin{enumerate}
\item 求$X$的数学期望$EX$(记$EX$为$b$)；
\item 求$\mu$的置信度为0.95的置信区间；
\item 利用上述结果求$b$的置信度为0.95的置信区间．
\end{enumerate}
\end{problem}


\makepart{}{本题满分8分}

\begin{problem}
设$A,B$是二随机事件；随机变量
$$X = \begin{cases}
   1, & \text{若$A$出现} \\
  -1, & \text{若$A$不出现}
\end{cases}\qquad Y = \begin{cases}
   1, & \text{若$B$出现} \\
  -1, & \text{若$B$不出现}
\end{cases}$$
试证明随机变量$X$和$Y$不相关的充分必要条件是$A$与$B$相互独立．
\end{problem}


\end{document}
